\scanpage{95}

It makes \unsure{our hands|our hand} itch to \deleted{\unsure{fool around}}\inserted{work}
in $\mathrm{Gl}(2,\mathbb Z)$ itself, a central extension of
$\mathfrak S_{0,3}$ --- in fact, this extension is interpreted as
$\mathcal{T}_{1,1}$ --- but we shall come back to it, later\ldots

\uline{V}) Let us sum up the avatars of the same group, which we have
\unsure{encountered|encountered} so far

1) $\mathfrak S_{0,3}\overset{\text{dfn}}{=}\operatorname{Aut}_{\mathrm{lac}}(\Pi_{0,3})$
with \unsure{subgroup} $\mathfrak S_{0,3}^{+}$ (preserving orientation)

2) $\mathcal{T}'_{0,4}\simeq\operatorname{Autext}_{\mathrm{lac}}(\Pi_{0,4})'$
outer automorphisms of $\Pi_{0,4}$ respecting the loop structure and fixing the loop class
$\unsure{L_4|L_\lambda}$ (and permuting $L_0,L_1,\unsure{L_\infty|L_3}$)

\deleted{with $\unsure{\mathcal{T}_{0,4}^{\prime+}}$}
\deleted{$\unsure{A'_{0,4}/A_{0,4}^{\prime0}}\simeq\unsure{\operatorname{Auttop}}(\mathbb P^1_{\mathbb C}\setminus\{0,\unsure{1}\})$}
with $\mathcal{T}_{0,4}^{\prime+}$ (preserving orientation)

3) $A'_{0,4}/A_{0,4}^{\prime0}$ group of connected components of
\[
A'_{0,4}\subset\operatorname{Auttop}
\left(\deleted{\mathbb P^1_{\mathbb C}\setminus\unsure{\{}}
\underset{\substack{\Vert\\\mathbb P^1_{\mathbb C}\setminus\{0,1,\infty;\tfrac12\}}}{U_{0,4}^{\mathrm{top}}}\right)
\]
subgroup of automorphisms that fix the point at infinity $a_{0,4}=\tfrac12$

with $A_{0,4}^{\prime+}/A_{0,4}^{\prime0}$ (preserving orientation

4) $\Pi_1\bigl((U_{0,3}^{\mathrm{an}},\overbrace{\mathfrak S_3\times\unsure{\mathbb Z_2|\mathbb Z/2}}^{\mathcal{T}_{0,3}}),-\bar j\bigr)$
(described in the manner of J. Malgoire)
with the subgroup $\Pi_1\bigl((U_{0,3}^{\mathrm{an}},\mathfrak S_3),-\bar j\bigr)$
$(\mathfrak S_3=\mathcal{T}_{0,3}^{+})$

5) $\Pi_1\bigl((\mathbb P^{1\,\mathrm{an}}_{\mathbb C},0\cdot\{0\}+2\{1\}+3\{\infty\}),-\bar j\bigr)$
(in the manner of J. Malgoire)
with the subgroup $\Pi_1\bigl(\mathbb P^{1\,\mathrm{an}}_{\mathbb C},\unsure{0}\{0\}+2\{1\}+3\{\infty\}\bigr)$

6) Triangulated cartographic group \deleted{\illegible{1}} $\Gamma_{\infty,3}$
\[
\Gamma_{\infty,3}\unsure{\}}=\{r_0,r_1,r_\infty\mid
r_0^2=r_1^2=r_\infty^2=(r_0r_\infty)^2=(r_1r_0)^3=1\}
\]
(with oriented triangulated cartographic subgroup
\[
\Gamma_{\infty,3}^{\circ}=\{
\underset{\substack{\Vert\\r_\infty r_1}}{\rho_s},
\underset{\substack{\Vert\\r_0r_\infty}}{\sigma},
\underset{\substack{\Vert\\r_1r_0}}{\rho_f}
\mid\rho_f\sigma\rho_s=\sigma^2=\rho_f^3=1\}.
\]
)

\margin{this is also the group
$\widetilde{\Pi}_{0,3}/(\ell_1^2,\ell_\infty^3)$, where
$\widetilde{\Pi}_{0,3}=\Pi_{0,3}\cdot\{1,\tau\}$}

7) $\Pi_1\bigl(\mathcal U_{0,3}^{\mathrm{conf}},\operatorname{pt}(U_{0,3},-\bar j)\bigr)$
\inserted{universal conformal surface of type $(0,3)$}
with subgroup $\Pi_1\bigl(\mathcal U_{0,3}^{\mathrm{an}},\operatorname{pt}(U_{0,3},-\bar j)\bigr)$
\inserted{universal holomorphic surface of type $(0,3)$}

\margin{\uline{N.B.} $\mathcal U_{0,3}^{\mathrm{conf}}\simeq(U_{0,3},\mathcal T_{0,3})$
which makes it possible to identify 7) with 4)}

8) $\mathcal T'_{1,1}=\mathcal T_{1,1}/\underbrace{\{\pm1\}}_{\text{center}}$,
with $\mathcal T_{1,1}\simeq\operatorname{Autext}_{\mathrm{lac}}(\Pi_{1,1})
\simeq A_{1,1}/A_{1,1}^{0}$

($A_{1,1}\simeq\text{topological automorphisms of }\unsure{U_{1,1}^{\circ}|U_{1,1}}$)
with subgroup $\mathcal T_{1,1}^{\prime+}=\mathcal T_{1,1}^{+}/\{\pm1\}$

\scanpage{96}
\mspage{337}

9) $\Pi_1(M_{1,1}^{\mathrm{conf}},\mathrm{pt})' \simeq \unsure{\Pi_1()/\mathrm{centre}\{\pm1\}} \simeq \Pi_1(M_{1,1}'{}^{\mathrm{conf}},\mathrm{pt})$

where $M_{1,1}^{\mathrm{an}}$ is the \deleted{analytic}\inserted{\unsure{conformal|conformal}} modular multiplicity of \inserted{conformal} elliptic curves (i.e. \deleted{analytic}\inserted{conformal} curves of type $1,1$), $M_{1,1}'{}^{\mathrm{an}}$ that of \deleted{analytic}\inserted{conformal} elliptic curves ``modulo symmetry''; with subgroup
\[
\Pi_1(M_{1,1}^{\mathrm{an}},\mathrm{pt})' \simeq \unsure{\Pi_1()/\mathrm{centre}\{\pm1\}} \simeq \Pi_1(M_{1,1}'{}^{\mathrm{an}},\mathrm{pt})
\]
where the base point is the \inserted{conformal elliptic \unsure{resp.}} analytic curve, quotient of $\mathbb C$ by $\mathbb Z+\mathbb Z(-\bar j)$. \uline{N.B.} The modular multiplicities are $(D,\mathrm{Gl}(2,\mathbb Z))$, $(D,\mathrm{Sl}(2,\mathbb Z))$ in the \deleted{analytic} conformal resp. analytic setting for $M_{1,1}$, and $(D,\mathrm{Gl}(2,\mathbb Z)'=\mathrm{GP}(1,\mathbb Z))$ resp. $(D,\unsure{\mathrm{Sl}(1,\mathbb Z)'|\mathrm{Sl}(2,\mathbb Z)'}=\mathrm{GP}(1,\mathbb Z)^+)$ in the conformal resp. analytic setting for $M_{1,1}'$.

10) $\mathrm{Gl}(2,\unsure{\mathbb Z^{\ast}|\mathbb Z})' \simeq \mathrm{GP}(1,\mathbb Z)$, with the subgroup $\mathrm{Sl}(2,\mathbb Z)'\simeq\mathrm{GP}(1,\mathbb Z)^+$.

\scanpage{97}

(VI) \uline{The loop automorphisms of} $\widehat{\Pi}_{0,3}$ \uline{and of} $\widehat{\mathfrak S}_{0,3}^{+}$

We have a canonical \unsure{commutative} diagram:

(49)
\diagram{\begin{tikzcd}[ampersand replacement=\&,column sep=large,row sep=large]
\begin{array}{l}
\operatorname{Autext}(\widehat{\mathfrak S}_{0,3}^{+})'\\
=\text{group of outer automorphisms}\\
\text{inducing}\\
\text{the identity in the}\\
\text{quotient }\mathfrak S_3\simeq\mathcal T_{0,3}^{+}
\end{array}
\arrow[r,"{\text{restriction}}","\sim"']
\&
\begin{array}{l}
\operatorname{Autext}^{\unsure{\rho}}(\widehat{\Pi}_{0,3})'\\
=\text{group of outer automorphisms}\\
\text{that}\\
\text{commute with}\\
\text{the outer action of }\mathfrak S_3
\end{array}
\\
\begin{array}{l}
\operatorname{Aut}(\widehat{\mathfrak S}_{0,3}^{+},\sigma_\infty)'\\
=\text{group of automorphisms}\\
\text{of }\widehat{\mathfrak S}_{0,3}^{+}\text{ that }\deleted{\unsure{\text{fix}}}\inserted{\text{fix}}\ \sigma_\infty,\\
\text{and that fix }\rho\ \mathrm{mod}\ \widehat{\Pi}_{0,3}
\end{array}
\arrow[u,"\sim"]\arrow[r,"{\text{restriction}}","\sim"']
\&
\begin{array}{l}
\operatorname{Aut}(\widehat{\Pi}_{0,3},\sigma_\infty)\\
=\text{group of automorphisms}\\
\text{that commute}\\
\text{with }\sigma_\infty,\text{ and that}\\
\text{commute}\\
\mathrm{ext}^{\mathrm t}.\text{ with }\rho
\end{array}
\arrow[u,"\sim"]
\end{tikzcd}}

\margin{\deleted{\illegible{1}}}

The \unsure{$\in$} \deleted{\unsure{elements}} $\operatorname{Aut}(\widehat{\mathfrak S}_{0,3}^{+},\sigma_\infty)'$ are of the form
\[
(50)\qquad
\left\{\begin{aligned}
u(\sigma_\infty)&=\sigma_\infty\\
u(\rho)&=g\rho
\end{aligned}\right.
\qquad (\text{i.e. }u\rho u^{-1}\rho^{-1}=g)
\]

\margin{We write the formulas in $\operatorname{Aut}(\widehat{\mathfrak S}_{0,3}^{+})'=G$, group of automorphisms of $\widehat{\mathfrak S}_{0,3}^{+}$ that induce the identity in $\widehat{\mathfrak S}_{0,3}^{+}/\Pi_{0,3}\simeq\mathfrak S_3$, \uline{\unsure{is identified}} with the group $\operatorname{Aut}(\Pi_{0,3})'=\mathrm{gr.}$ of automorphisms of $\Pi_{0,3}$ that commute \inserted{\unsure{\uline{outer}}} with $\mathfrak S_3$; -- we shall have $\Pi_{0,3}\subset\widehat{\mathfrak S}_{0,3}^{+}\subset\widehat{\mathfrak S}_{0,3}\subset\unsure{\mathfrak g.|g.}\,G$. \uline{N.B.} We assume here that the center of $\widehat{\mathfrak S}_{0,3}^{+}$ is trivial.}

$g\in\widehat{\Pi}_{0,3}$ \deleted{\illegible{1}} is subject to the cocycle condition
\[
(51)\qquad \boxed{g\rho(g)\rho^2(g)=1}
\]
which expresses the relation
\[
(52)\qquad (g\rho)^3=1\ .
\]

We thus find a $1-1$ correspondence between the elements of $\operatorname{Aut}(\widehat{\mathfrak S}_{0,3}^{+},\sigma_\infty)'$ and the $g\in\widehat{\Pi}_{0,3}$ \unsure{satisfying} (51), and such moreover that

(53) $\sigma_\infty,g\rho$ with the relations $\sigma_\infty^2=(g\rho)^3=1$ \unsure{present} $\widehat{\mathfrak S}_{0,3}^{+}$

\scanpage{98}
\mspage{338}

(which I am going to call the ``preliminary condition'').

\deleted{For \unsure{$u$ thus defined by a}}

For a $u\in G=\operatorname{Aut}(\widehat{\mathfrak S}_{0,3}^{+})'\unsure{!}$ (not necessarily fixing $\sigma_\infty$), we have equivalent conditions
\[
(54)\quad
\left\{\begin{aligned}
\text{(i)}\quad &u(\overset{\inserted{\varepsilon_0}}{\ell'_0})\in\text{conjugate of }\deleted{\unsure{\ell}}\varepsilon_0^{\widehat{\mathbb Z}}=L'_0\ (\text{in }\widehat{\mathfrak S}_{0,3}^{+})\\
\text{(ii)}\quad &u(\ell_0)\in\text{conjugate of }\ell_0^{\widehat{\mathbb Z}}=L_0\ (\text{in }\widehat{\Pi}_{0,3})\\
\text{(iii)}\quad &u(L_i)\text{ conjugate of }L_i\quad\text{for }i=0,1,\infty\\
\text{(iv)}\quad &\exists\ p\in\widehat{\mathbb Z}^{\ast},\ g_0\in\widehat{\Pi}_{0,3}\quad
\boxed{u(\varepsilon_0)=\operatorname{int}(g_0)\varepsilon_0^p}\\
\text{(v)}\quad &\exists\ p\in\widehat{\mathbb Z}^{\ast},\ g_0\in\widehat{\Pi}_{0,3}\quad
\boxed{u(\ell_0)=\operatorname{int}(g_0)(\ell_0^p)}
\end{aligned}\right.
\]

We shall then also have
\[
(55)\qquad
\begin{aligned}
\deleted{\unsure{u}}\quad u(\ell_1)&=\operatorname{int}(g_1)(\ell_1^p)&
g_1&=g\rho(g_0)\\
u(\ell_\infty)&=\operatorname{int}(g_\infty)(\ell_\infty^p)&
g_\infty&=g\rho(g_1)=g\rho(g)\rho^2(g_0)\ .
\end{aligned}
\]

The element \deleted{$p\in\widehat{\mathbb Z}^{\ast}$ \unsure{is}}
\[
(55)\qquad p\in\widehat{\mathbb Z}^{\ast}\qquad p=\chi(u)
\]
is well determined by $u$, hence by $g$\inserted{(if $u$ given by (50))}, we shall call it the \uline{multiplier} of $u$. The element \deleted{\unsure{$g_0$}}
\[
(56)\qquad g_0\in\widehat{\Pi}_{0,3}
\]
is determined a priori modulo $L_0=\ell_0^{\widehat{\mathbb Z}}$ (change
\[
(57)\qquad g_0\longmapsto g_0\ell_0^\alpha
\]
).

\margin{Notations $\operatorname{Autext}_{\mathrm{lac}}(\widehat{\mathfrak S}_{0,3}^{+})$, $\operatorname{Aut}_{\mathrm{lac}}(\widehat{\mathfrak S}_{0,3}^{+})$, $\operatorname{Aut}_{\mathrm{lac}}(\widehat{\mathfrak S}_{0,3}^{+})'$, $\operatorname{Aut}_{\mathrm{lac}}(\widehat{\mathfrak S}_{0,3}^{+},\sigma_\infty)'$, $\operatorname{Aut}_{\mathrm{lac}}(\widehat{\Pi}_{0,3},\sigma_\infty)\ldots$}

If the \inserted{equivalent} conditions (\unsure{54|55}) are satisfied, we say that $u$ respects the loop structure of $\widehat{\mathfrak S}_{0,3}^{+}$ (or of $\Pi_{0,3}$) -- it being understood that $\widehat{\mathfrak S}_{0,3}^{+}$ is equipped with the ``loop-group'' $\varepsilon_0^{\widehat{\mathbb Z}}\simeq L'_0$. We \unsure{shall restrict|restrict} ourselves from now on to such automorphisms. \uline{N.B.} These conditions depend only on the outer automorphism defined by $u$, in $\widehat{\Pi}_{0,3}$ or in $\widehat{\mathfrak S}_{0,3}^{+}$. We

\scanpage{99}

are going to \unsure{take an interest} in the canonical liftings of such outer automorphisms given by (49), (50), to actual automorphisms.

The defining formula (50) for $u$ involves only $g$\inserted{$\in\widehat{\Pi}_{0,3}$}, in principle $g_0\ \mathrm{mod}\ L_0$ in $\widehat{\Pi}_{0,3}/L_0$ is deduced from $g$. But in practice $g_0$ plays as important a role as $g$ -- we shall see that to a large extent it determines $g$ (formula (62)).

If there were no need to introduce $g_0$, to express the conditions (54), we could easily \uline{parametrize} the $g$ \unsure{satisfying} (51) -- hence corresponding to \deleted{\unsure{splittings of}}\inserted{liftings of} $\dot\rho\in\mathfrak S_3$ to elements of order $3$ of $\widehat{\mathfrak S}_{0,3}^{+}$ -- by noting that there are exactly two conjugacy classes of such liftings, (corresponding to the fixed points $-\bar j,-j$ of $\dot\rho$ operating on $U_{0,3}^{\mathrm{an}}$), which are parametrized by
\[
(58)\quad
\left\{\begin{aligned}
g&=\lambda_0^{\unsure{\sharp|\ast}}\rho(\lambda_0^{\unsure{\sharp|\ast}})^{-1}
=[\lambda_0^{\unsure{\sharp|\ast}},\rho]
&&\begin{gathered}\text{(principal \deleted{first}\inserted{class}, \unsure{corresponding|corresponding}}\\
\text{to the ``north pole'' }-\bar j)\end{gathered}\\
g&=\lambda_0^{\unsure{\sharp|\ast}\prime}\ell_1^{-1}\rho(\lambda_0^{\unsure{\sharp|\ast}\prime})^{-1}
&&\begin{gathered}\text{(adjoint class, corresponding to the}\\
\text{``south pole'' }-j)\end{gathered}
\end{aligned}\right.
\]

\margin{\unsure{These|The} parametric \unsure{representations} are \unsure{unambiguous|unique}.}

where $\lambda_0,\lambda_0^{\unsure{\sharp|\ast}\prime}\in\widehat{\Pi}_{0,3}$. We shall see below that these cases correspond respectively to the cases
\[
(59)\qquad p\equiv1\ (3)\qquad\text{and}\qquad p\equiv-1\ (3)\ .
\]